\section{Постановка проблемы}
\label{sec:bgp-dynamics}

\subsection{Нестабильность и пересылка избыточных обновлений}

Серия работ Лабовица и соавторов на рубеже 2000-х показала, что
междоменная маршрутизация в реальной сети ведёт себя совсем не как
статическая структура. В \textit{Internet Routing
Instability}~\cite{labovitz1997instability} зафиксировано, что объём
BGP-обновлений в среднем кратно превосходит интуитивный сценарий
«редкие, осмысленные изменения». Существенная доля сообщений
повторная, избыточная или возникает из-за патологий смежных
маршрутизаторов. Любой прикладной сервис, привязанный к
BGP-состоянию, обязан рассчитывать на непрерывные изменения как на штатный
режим, а не как на аномалию.

В \textit{Delayed Internet Routing Convergence}
\cite{labovitz2000convergence} та же картина дополнена временным
срезом: после изменения топологии глобальная сходимость занимает
минуты, а в худших сценариях -- значительно больше. Для прикладной
системы, опирающейся на отображение
\textit{префикс~$\rightarrow$~ASN}, неопределённость возникает на двух
уровнях:

\begin{enumerate}
\item собственно динамика BGP, заданная распределённой природой
протокола и не поддающаяся одностороннему контролю;
\item дополнительная задержка локального обновления таблицы внутри
самой прикладной системы.
\end{enumerate}

Первое нельзя устранить -- его остаётся только принять как свойство
среды. Второе целиком определяется архитектурой сервиса, и именно его
снижением занимается настоящая работа.

\subsection{Задача устойчивых путей}

Гриффин, Шеферд и Вилфонг~\cite{griffin2002stablepaths} формализовали
работу BGP как Stable~Paths~Problem~(SPP). Модель показывает, что
протокол ищет локально устойчивый выбор маршрутов в присутствии
политик, и этот процесс в общем случае не обязан сходиться к единому
глобальному оптимуму. Для прикладной части отсюда следуют два
наблюдения.

Во-первых, у BGP нет «идеальной глобальной истины»: один и тот же
префикс одновременно может выглядеть по-разному с точки зрения разных
наблюдателей в зависимости от их положения в графе и от их собственных
политик.

Во-вторых, если даже сам протокол достигает согласия только в пределе,
то нет смысла строить прикладной сервис как редкий синхронный
пересчёт. Гораздо ближе к реальности здесь модель, в которой состояние
эволюционирует во времени и обновляется отдельными событиями.

\subsection{Роль публичных коллекторов маршрутов}

Платформы RIPE~Routing~Information~Service~\cite{ripe-ris} и
RouteViews~\cite{routeviews} многие годы служат основными источниками
данных как для академических исследований BGP, так и для
эксплуатационной аналитики. RIS опирается на распределённую сеть
коллекторов, собирающих маршрутные данные у добровольных партнёров на
крупных точках обмена трафиком и через multihop-сеансы. RouteViews
предоставляет доступ к текущим и историческим данным коллекторов и к
архивам в формате MRT.

Для настоящей работы важно, что именно эти платформы давно живут в
двойной модели данных: периодические дампы таблиц плюс непрерывный
поток обновлений. Иначе говоря, даже на исследовательском уровне
схема «снимок + поток» считается естественным способом организации
маршрутных данных. Хорошая иллюстрация -- механизм RIS Beaconing
\cite{ripe-beacons}, в котором заранее известные префиксы объявляются
и отзываются по расписанию, превращая BGP в контролируемый
эксперимент.

\subsection{Ограничения пакетной модели обновления}

Допустим, прикладная система раз в несколько часов целиком
перестраивает таблицу \textit{префикс~$\rightarrow$~ASN} по BGP
полной таблице маршрутов. У такого подхода есть понятные преимущества:
не нужно
отслеживать инкрементальные события, не нужно следить за порядком
обновлений, а внутреннее хранилище можно оптимизировать под чтение и
под редкую целиковую замену.

В динамической среде BGP такая схема упирается сразу в несколько
ограничений:

\begin{enumerate}
\item между фактическим изменением маршрута и его появлением в
локальной таблице возникает лаг; для систем реального времени лаг в
часы -- это уже работа с устаревшей картиной мира;
\item чем чаще приходится перестраивать таблицу, тем больше времени
занимает поддержание состояния: время одной перестройки определяется размером
таблицы $n$, а не объёмом фактических изменений;
\item локальная правка одного префикса требует работы, сопоставимой
по объёму с полной сборкой таблицы заново;
\item с ростом числа маршрутов (сегодня IPv4 уже превысил миллион, а
IPv6 -- сотни тысяч префиксов) и частоты их изменений пакетная
схема начинает не укладываться в требования по задержке и по
потреблению CPU.
\end{enumerate}

Отсюда естественный переход к другой схеме: начальный снимок нужен
только при старте, а дальше актуальность таблицы поддерживается
потоком обновлений.

\subsection{Локальная динамика и flap damping}

Динамика BGP не сводится к «полезным» изменениям топологии. Заметная
часть update-трафика порождается локальными нестабильностями --
короткоживущими событиями, повторными withdraw/announce одного и того
же префикса, последствиями переключения на резервные пути. RFC~2439
\cite{rfc2439} описывает классический механизм route~flap~damping,
который подавляет излишне нестабильные маршруты экспоненциально
убывающим штрафом.

Прикладному сервису flap damping не нужно реализовывать
самостоятельно, но он остаётся важной иллюстрацией двух более общих
наблюдений:

\begin{itemize}
\item поток BGP-обновлений несёт значительную долю избыточной
информации;
\item потоковая часть сервиса должна выдерживать всплески,
многократно превышающие средний поток событий.
\end{itemize}

Это задаёт требование к структуре хранения: операция обновления должна
изменять затронутый префикс или шард, а не зависеть от полного размера
таблицы $n$. В экспериментальной части это требование проверяется по
времени \texttt{Upsert}/\texttt{Delete}, объёму выделенной памяти и числу
аллокаций. Вместе с семантикой \texttt{Lookup} оно определяет выбор
алгоритмов в главах~\ref{sec:problem} и~\ref{sec:art-storage}.

\subsection{Системные следствия для прикладного сервиса}

Из рассмотренных свойств BGP следуют четыре требования к прикладному
сервису.

\begin{enumerate}
\item Задержка актуализации таблицы должна определяться обработкой
отдельных BGP-событий, а не периодом редкого полного пересчёта.
\item Рассинхронизация неизбежна; правильное проектирование закладывает
дисциплину пересинхронизации как штатный режим, а не как обработку
аварии.
\item Время обработки одного события на стороне сервиса должно зависеть
от локального изменения состояния, иначе потоковая модель сводится к
повторной полной перестройке.
\item Семантика \texttt{Lookup} не должна зависеть от того, в какой
момент применилось конкретное обновление: видимое клиенту состояние
обязано соответствовать какой-то определённой точке последовательности
применённых событий.
\end{enumerate}

Эти четыре положения становятся техническими требованиями к сервису,
формализованными в главе~\ref{sec:problem} как функциональные и
нефункциональные ограничения.

%==============================================================
